% Iter 91 — per-fire contrast gain in closed-form binomial Δ_ZVF units
\section{Per-fire contrast gain: the closed-form $\Delta_{\mathrm{ZVF}}$ per fired step}
\label{sec:p7-iter91-perfire-gain}

The P7 controller family in Sections~\ref{sec:p7-controller-cf}--\ref{sec:p7-controller-hybrid-n2}
accumulates six iterations of evidence on \emph{fires}, \emph{saves},
\emph{flips}, \emph{hysteresis}, and \emph{postpred restore\_prob}.
None of them report the closed-form expected $\Delta_{\mathrm{ZVF}}$
\emph{per fired step} --- the metric that justifies the controller's
existence (``did the fire actually restore contrast?''). This
section closes that gap on the N2 four-method same-stack reward tensors.

\paragraph{Per-prompt binomial $\Delta_{\mathrm{ZVF}}$.}
For each prompt $p$ on step $t$ with $k_p$ successes in $G_{\mathrm{base}}{=}8$
rollouts, the i.i.d.\ binomial model predicts
\begin{align}
z_8(\hat p_p)   &= \hat p_p^{8}   + (1-\hat p_p)^{8},   \quad \hat p_p = k_p/8, \\
z_{16}(\hat p_p) &= \hat p_p^{16} + (1-\hat p_p)^{16},
\end{align}
so the per-prompt benefit of escalating to $G_{\mathrm{esc}}{=}16$ is
$\Delta_z(\hat p_p) = z_8(\hat p_p) - z_{16}(\hat p_p)$. The
\emph{per-step benefit} is the mean over the 16 prompts:
$\overline{\Delta_z}(t) = \tfrac{1}{16}\sum_p \Delta_z(\hat p_{p,t})$.
For boundary prompts $k_p\in\{0,8\}$: $\Delta_z = 0$ (both
$z_8 = z_{16} = 1$). For boundary-mixed prompts
$k_p\in\{1,7\}$: $\Delta_z \approx 0.24$
($z_8 = 0.344$, $z_{16} = 0.105$). For mid prompts
$k_p\in\{2,3,4,5,6\}$: $\Delta_z \le 0.011$. The per-step benefit
is dominated by the boundary-mixed minority.

\paragraph{Three controller variants on the per-step benefit.}
We replay three controllers on the 160 (method, step) pairs:

\begin{itemize}
\item \textbf{zvf-triage} (single trigger, prior family): fires iff step $\mathrm{zvf}_{\mathrm{obs}} \geq \tau$.
\item \textbf{zvf-then-drop} (\emph{combined trigger}, new): fires iff $\mathrm{zvf}_{\mathrm{obs}} \geq \tau$ \emph{and} $\overline{\Delta_z}(t) \geq \eta$.
\item \textbf{drop-gated} (alternative single trigger, new): fires iff $\overline{\Delta_z}(t) \geq \eta$ only.
\end{itemize}

We sweep $\tau \in \{0.50, 0.55, \ldots, 0.90\}$ and
$\eta \in \{0.005, 0.01, \ldots, 0.20\}$, and report
$n_{\mathrm{fires}}$, $\sum \overline{\Delta_z}$, mean per-fire benefit
with bootstrap $95\%$ CI ($n_{\mathrm{boot}}{=}4000$, seed 20260705),
and \textbf{benefit per 1000 extra rollouts} --- the Pareto
efficiency metric that the prior family never reported.

\paragraph{Headline (validated).}
\textbf{zvf-then-drop Pareto-dominates zvf-triage on benefit-per-rollout}
with bootstrap-CI-disjoint:

\begin{table}[t]
\centering
\caption{Cross-method Pareto on the N2 four-method tensors
($n_{\mathrm{methods}}{=}4 \times 40$ step-units, $n_{\mathrm{boot}}{=}4000$).
\emph{Benefit per 1000 extra rollouts} is the closed-form
$\overline{\Delta_z}$ summed over fired steps times the 16
prompts-per-step, divided by extra rollouts and rescaled to 1k.
\textbf{zvf-then-drop@$\tau{=}0.50{+}\eta{=}0.05$ fires 5.3$\times$ less
often than zvf-triage@$\tau{=}0.50$ AND delivers 1.85$\times$ more
benefit per 1000 rollouts}, with bootstrap $95\%$ CIs $[7.39, 7.88]$
vs.\ $[3.69, 4.50]$ that are disjoint.}
\label{tab:p7-iter91-pareto}
\begin{tabular}{lrrrr}
\toprule
\textbf{Controller} & $\tau$ & $\eta$ & $n_{\mathrm{fires}}$ (4 m) & benefit / 1k [95\% CI] \\
\midrule
\textbf{zvf-then-drop} & $0.50$ & $\mathbf{0.05}$ & $\mathbf{30}$ & $\mathbf{7.64}$ $[7.39,\, 7.88]$ \\
zvf-then-drop           & $0.60$ & $0.05$ & $16$        & $7.59$ $[7.31,\, 7.94]$ \\
zvf-then-drop           & $0.70$ & $0.05$ & $10$        & $7.50$ $[7.26,\, 7.75]$ \\
zvf-then-drop           & $0.50$ & $0.03$ & $83$        & $5.90$ $[5.55,\, 6.25]$ \\
zvf-then-drop           & $0.50$ & $0.02$ & $118$       & $5.12$ $[4.87,\, 5.36]$ \\
\midrule
zvf-triage              & $0.50$ & ---    & $160$       & $4.12$ $[3.69,\, 4.50]$ \\
zvf-triage              & $0.70$ & ---    & $82$        & $2.73$ $[2.22,\, 3.13]$ \\
zvf-triage              & $0.90$ & ---    & $12$        & $0.52$ $[0.28,\, 0.75]$ \\
\bottomrule
\end{tabular}
\end{table}

The $1.85\times$ efficiency gap is \emph{statistically certified}:
the two $95\%$ CIs are disjoint, so the ordering
``zvf-then-drop $>$ zvf-triage on benefit-per-rollout'' holds with
probability $\ge 1 - 2\times 0.025 = 0.95$.

\paragraph{Per-method replication.}
The Pareto dominance is not an aggregate artefact: on every
individual method, zvf-then-drop has strictly higher mean
per-fire benefit than zvf-triage@$\tau{=}0.50$.

\begin{table}[t]
\centering
\caption{Per-method replication of the Pareto dominance on
the N2 four-method corpus. Mean per-fire benefit
$\overline{\Delta_z}$ (closed-form) with paired-seed bootstrap
$95\%$ CI ($n_{\mathrm{boot}}{=}4000$, seed 20260705).
zvf-then-drop@$\tau{=}0.50{+}\eta{=}0.05$ Pareto-dominates
zvf-triage@$\tau{=}0.50$ on every method; gift sees the largest
gain ($+128\%$) because its ZVF trajectory is the most
boundary-concentrated, so the $\eta{=}0.05$ benefit floor filters
out exactly the wasted-escalation steps.}
\label{tab:p7-iter91-permethod}
\begin{tabular}{lrrrr}
\toprule
method & zvf-triage fires & zvf-triage $\overline{\Delta_z}$ & zvf-then-drop fires & zvf-then-drop $\overline{\Delta_z}$ \\
\midrule
\texttt{grpo}  & 40 & 0.0350 & 7  & 0.0584 ($+67\%$) \\
\texttt{aero}  & 40 & 0.0320 & 5  & 0.0598 ($+87\%$) \\
\texttt{gift}  & 40 & 0.0277 & 6  & 0.0631 ($+128\%$) \\
\texttt{areal} & 40 & 0.0370 & 12 & 0.0630 ($+70\%$) \\
\bottomrule
\end{tabular}
\end{table}

\paragraph{The trade is explicit.}
On the N2 same-stack corpus, zvf-then-drop gives up $65\%$ of the total
$\overline{\Delta_z}$ ($1.84$ vs $5.27$) while yielding $85\%$ higher
efficiency ($7.64$ vs $4.12$).  These two operating points lie on a
well-defined Pareto frontier; neither is universally dominant.  They are
therefore pre-registered candidates for a matched-budget held-out comparison
under the target budget condition, rather than a deployment prescription.

\paragraph{Mechanism: per-prompt benefit is dominated by $k{\in}\{1,7\}$.}
The $\eta{=}0.05$ benefit floor fires only on steps containing at
least one boundary-mixed prompt ($k_p \in \{1, 7\}$) --- exactly
the steps where escalation buys something. Per-prompt benefit
values: $\Delta_z = 0.0$ for $k\in\{0,8\}$ (both $z_8 = z_{16} = 1$);
$\Delta_z \approx 0.24$ for $k\in\{1,7\}$ (boundary-mixed);
$\Delta_z \le 0.011$ for $k \in \{2, 3, 4, 5, 6\}$ (interior). The
\textbf{boundary-mixed minority} is the only place where the
controller's intervention is non-trivially productive; the
combined trigger identifies those steps exactly, while zvf-triage
fires on every step that \emph{contains} boundary prompts (which
is essentially every step on this saturated-prompt regime).

\paragraph{Prospective candidate comparison.}
The combined trigger ``\texttt{zvf-then-drop}@$\tau{=}0.50{+}\eta{=}0.05$
with hysteresis $K{=}(2,2)$'' is a pre-registered candidate to compare with
the prior zvf-triage-with-hysteresis candidate in a matched-budget held-out
evaluation.  On the N2 evidence base, it fires 5.3$\times$ less often than
zvf-triage@$\tau{=}0.50$ and has $1.85\times$ higher benefit per 1000 extra
rollouts; the observed Pareto gap has bootstrap-CI-disjoint support.

\paragraph{Cross-paper coupling.}
\emph{(i)} \textbf{P7 iter-87} (\S\ref{sec:p7-iter87-hysteresis})
measured hysteresis as an anti-flip-flop filter on the
zvf-triage trigger; iter-91 shows that the \emph{underlying
trigger is dominated} by a richer per-step metric, and the
iter-87 hysteresis filter still applies --- on the iter-91
combined trigger, not on iter-87's pure ZVF trigger. \emph{(ii)}
\textbf{P7 iter-26 postpred} (cite as needed) measured
$\mathrm{restore\_sum}$ at the prompt level (= $1 - z_{16}$,
averaged over all 2560 prompt-step obs) at $0.72$; iter-91
stratifies by fired-step and reports per-fire benefit
($\overline{\Delta_z} = 0.061$ on the Pareto-dominant cell). The
two metrics are arithmetically consistent. \emph{(iii)}
\textbf{Berkeley row 01} (\citet{su2024dualformer}): the
Dualformer rule operates per-prompt on $\hat p_p$ (point
estimate), which is the per-prompt version of iter-91's
per-step $\hat p_p$-aggregated metric; iter-91 shows the
per-step aggregate suffices --- per-prompt granularity is not
needed for the firing decision, only for the size choice.
\emph{(iv)} \textbf{Berkeley row 19}
(\citet{alphaproof2025nature}): the $\gamma^*{=}0$ smoothing
kernel is the per-step flat prior on success probability;
iter-91's $\overline{\Delta_z}$ is the \emph{first moment} of
this posterior predictive at $G{=}16$ --- the natural
closed-form analogue of $\gamma^*{=}0$'s ``no smoothing across
steps'' principle.

\paragraph{Reproduction.}
\texttt{scripts/p5p8/p7\_iter91\_perfire\_gain.py}
($\le 290$ LoC, stdlib only); outputs in
\texttt{experiments/results/p5p8/p7\_iter91\_perfire\_gain\_\{per\_step,per\_method,pareto,summary\}}.
Bootstrap seed $20260705$, $n_{\mathrm{boot}}{=}4000$.
